\documentclass[11pt,a4paper]{article}
\usepackage[utf8]{inputenc}
\usepackage[T1]{fontenc}
\usepackage[brazilian]{babel}
\usepackage[margin=2cm]{geometry}
\usepackage{lmodern}
\usepackage{amssymb}
\usepackage{xcolor}
\usepackage{array}
\usepackage{booktabs}
\usepackage{enumitem}
\usepackage{titlesec}
\usepackage{fancyhdr}
\usepackage{parskip}

\definecolor{brandgreen}{HTML}{2E5E3E}
\definecolor{lightbg}{HTML}{F5F2EC}

\titleformat{\section}{\large\bfseries\color{brandgreen}}{}{0pt}{}[{\color{brandgreen}\titlerule}]
\titleformat{\subsection}{\bfseries\color{brandgreen}}{}{0pt}{}

\pagestyle{fancy}
\fancyhf{}
\fancyhead[L]{\small Querência ERP --- Avaliação de usabilidade}
\fancyhead[R]{\small Fluxo 1 --- Associado}
\fancyfoot[C]{\small\thepage}
\renewcommand{\headrulewidth}{0.4pt}

\newcommand{\dizer}[1]{%
  \noindent\colorbox{lightbg}{%
    \begin{minipage}{0.95\linewidth}
    \vspace{1mm}
    \textbf{O que dizer ao participante:} ``#1''
    \vspace{1mm}
    \end{minipage}%
  }\vspace{3mm}
}
\newcommand{\criterio}[1]{\noindent\textit{Critério de sucesso objetivo:} #1}
\newcommand{\estado}[1]{\noindent\textit{Estado necessário (já preparado no kit):} #1\vspace{2mm}}

\begin{document}

\begin{center}
  {\LARGE\bfseries\color{brandgreen} Fluxo 1 --- Cadastro de associado}\\[1mm]
  {\large RF01, RF18 --- Papel: Associado (app mobile)}
\end{center}

\vspace{2mm}
\noindent Fazer \textbf{uma} das duas variantes abaixo (combinada previamente com o moderador).

\section*{Variante A --- Auto-cadastro}

\estado{ficha de persona ainda não usada.}

\dizer{Você quer se tornar associado da entidade e está criando sua conta pelo aplicativo pela
primeira vez.}

\begin{enumerate}[leftmargin=6mm]
  \item Partir da tela de login (app deslogado).
  \item Tocar \textbf{``Ainda não é associado? Cadastre-se''}.
  \item Preencher com os dados da ficha de persona do kit: Nome, CPF, Telefone (seção ``Dados
    pessoais''), Endereço (seção com ícone de mapa), E-mail e Senha.
  \item Tocar \textbf{Criar conta}.
  \item Esperado: mensagem de sucesso / tela indicando ``pendente de validação''.
\end{enumerate}

\textbf{Verificação objetiva:}
\begin{verbatim}
cd backend
node scripts/verificar-tarefas.js 1 <CPF-DA-FICHA>
\end{verbatim}
\criterio{associado existe com status \texttt{pendente\_validacao} \textbf{e} conta de login
vinculada.}

\section*{Variante B --- Vincular conta}

\estado{CPF pré-cadastrado pelo moderador (ativo, sem login), anotado no cartão do kit.}

\dizer{A diretoria já cadastrou você como associado, mas você nunca entrou no aplicativo. Use o
CPF que consta neste cartão para criar seu acesso.}

\begin{enumerate}[leftmargin=6mm]
  \item Partir da tela de login (app deslogado).
  \item Tocar \textbf{``Já é associado mas nunca acessou?''}.
  \item Preencher o CPF do cartão, e-mail e senha novos.
  \item Tocar \textbf{Vincular conta}.
\end{enumerate}

\textbf{Verificação objetiva:}
\begin{verbatim}
node scripts/verificar-tarefas.js 1 <CPF-DO-CARTAO>
\end{verbatim}
\criterio{ambos os critérios (status \texttt{ativo} + conta de login vinculada).}

\vspace{2mm}
\noindent\textit{Sem encadeamento: a conta criada neste fluxo não é usada nos Fluxos 2 e 3. O
moderador sai dela e entra na conta de teste do kit antes de seguir.}

\vfill
\noindent\footnotesize Se a tarefa travar por mais de \textasciitilde2 min, aplicar 1 intervenção
pontual e registrar; se mesmo assim não avançar, parar e marcar como \textit{não concluída} ---
protocolo completo no \texttt{roteiro-moderador.pdf}.

\end{document}
